% !TeX root = ../../Book3.tex
\section{Weierstrass 典范因子}
\label{b3:sec:2002}

直接相乘 \(1-z/a_n\) 往往失败，因为它的对数从一次项开始。乘上一个指数，把若干低阶项扣掉，尾部就可以变成更高次的可和项。

\subsection{初等因子}
\label{b3:subsec:200201}

定义\term[weierstrass-chudeng-yinzi]{Weierstrass 初等因子}[Weierstrass elementary factor]
\[
E_0(w)=1-w,\quad E_p(w)=(1-w)\exp\left(w+\frac{w^2}2+\cdots+\frac{w^p}p\right),\ p\ge1.
\]
它只有一个单零点 \(w=1\)，且 \(E_p(0)=1\)。指数部分不添加零点。

\subsection{典范因子}
\label{b3:subsec:200202}

当 \(|w|<1\)，固定对数分支给
\[
\log E_p(w)=-\sum_{k=p+1}^{\infty}\frac{w^k}k.
\]
所以 \(|w|\le1/2\) 时 \(|\log E_p(w)|\le2|w|^{p+1}\)，并可由 \(|e^v-1|\le e^{|v|}|v|\) 控制因子偏离一。若 \(\sum|a_n|^{-p-1}<\infty\)，则 \(\prod E_p(z/a_n)\) 在每个紧集正规收敛。这里 \(p\) 固定，而紧集半径只进入有限常数。

\subsection{属}
\label{b3:subsec:200203}

能使上述幂和收敛的最小非负整数 \(p\)，称该典范乘积的\term[dianfan-chengji-shu]{属}[genus]。当收敛指数 \(\lambda\) 不是整数时，\(p=\lfloor\lambda\rfloor\)；整数指数时，还要检查临界幂和是否收敛，不能只根据 \(\lambda\) 的值决定。

例如 \(a_n=n\) 的指数一，\(\sum1/n\) 发散而 \(\sum1/n^2\) 收敛，须用属一因子。若把正负整数成对相乘，一次指数正负抵消，可写成 \(\prod(1-z^2/n^2)\)；这依赖明确的配对，而不是许可随意重排一个条件收敛乘积。


初等因子的指数修正还控制对数导数。在避开 \(a_n\) 的地方，直接微分得到
\[
\frac{\mathrm{d}}{\mathrm{d}z}\log E_p(z/a_n)
=-\frac1{a_n-z}+\sum_{k=1}^{p}\frac{z^{k-1}}{a_n^k}
=-\frac{z^p}{a_n^p(a_n-z)}.
\]
当 \(|z|\le R\)、\(|a_n|>2R\) 时，最后一式的模不超过 \(2R^p/|a_n|^{p+1}\)。所以同一幂和不仅保证乘积存在，也保证尾部对数导数正规收敛。零点附近的留数由有限头部贡献，值为一；这将乘法的零点数据重新转成加法的主要部分数据。

\ProseParagraphBreak

对 \(a_n=n\)，属一乘积
\[
P(z)=\prod_{n\ge1}(1-z/n)e^{z/n}
\]
在原点满足 \(P(0)=1\)、\(P'(0)=0\)，且对数展开为
\[
\log P(z)=-\sum_{k\ge2}\frac{z^k}{k}\sum_{n\ge1}n^{-k},\quad |z|<1.
\]
换序由 \(\sum_{n,k}|z|^k/(kn^k)<\infty\) 保证。尤其二阶系数取决于 \(\sum n^{-2}\)，一次系数却因修正消失。这是初等因子“扣除低阶项”的具体可计算效果。

\ProseParagraphBreak

属在临界指数处的差别也能构造出来。点列 \(a_n=n\log^2(n+1)\) 的计数指数仍为一，但 \(\sum1/|a_n|\) 收敛，所以属零已够。点列 \(a_n=n\log(n+1)\) 的指数同样为一，而一次幂和发散，须用属一。二者的收敛指数相同，最小可用属不同。积分判别中出现的对数因子因此不能在临界问题中随意忽略。


\ProseParagraphBreak

若某个 \(a\) 应有重数 \(m\)，可以将它在点列中重复 \(m\) 次，也可以直接取因子 \(E_p(z/a)^m\)。对正规收敛的检验，后一写法须带上权重 \(m\)：充分条件成为 \(\sum m_n|a_n|^{-p-1}<\infty\)。遗漏重数会改变计数函数和收敛指数。例如在每个正整数 \(n\) 放重数 \(n\) 的零点时，带重数的幂和是 \(\sum n^{1-s}\)，收敛指数为二，属零或属一均不足；属二给 \(\sum n^{-2}\) 的尾和。

\ProseParagraphBreak

关于最小属还有一个容易混淆的区分。绝对收敛判据针对逐个初等因子的乘积；把具有对称性的因子成组后，低次项可能彼此抵消。例如正负整数的属一因子满足
\[
E_1(z/n)E_1(-z/n)=1-z^2/n^2.
\]
因此成对乘积有一个看似属零的二次因子表达，但它并没有使 \(\sum_{n\ne0}|n|^{-1}\) 收敛。解释“属”时必须说清楚使用单个零点的标准初等因子，还是已经经过配对的表达。后一表达式有利于计算，却不能反推前一幂和判据失效。
